% Separate proof increment. No input or consolidated checkpoint is modified.
% Unique label prefix: inc:AN03:local:v1:
\documentclass[11pt]{article}
\usepackage[T1]{fontenc}
\usepackage[utf8]{inputenc}
\usepackage{lmodern}
\usepackage[margin=0.9in]{geometry}
\usepackage{amsmath,amssymb,amsthm,mathtools,microtype,booktabs,array,xurl}
\usepackage[hidelinks]{hyperref}
\usepackage{fancyhdr}
\pagestyle{fancy}\fancyhf{}
\fancyhead[L]{\small AN03: local passage and mean tracking}
\fancyhead[R]{\small Review increment v1}
\fancyfoot[C]{\thepage}\setlength{\headheight}{14pt}
\setlength{\emergencystretch}{3em}\allowdisplaybreaks[2]
\numberwithin{equation}{section}
\newtheorem{theorem}{Theorem}[section]
\newtheorem{lemma}[theorem]{Lemma}
\newtheorem{proposition}[theorem]{Proposition}
\newtheorem{corollary}[theorem]{Corollary}
\theoremstyle{remark}\newtheorem{remark}[theorem]{Remark}
\newcommand{\dd}{\,d}\newcommand{\E}{\mathbb E}
\newcommand{\R}{\mathbb R}\newcommand{\esssup}{\operatorname*{ess\,sup}}
\title{Local-rate passage with collective mean tracking\\
\large A fixed-core theorem with a weighted tail-forcing budget}
\author{Analytical increment for independent review}
\date{Version 1 --- October 3, 2026}
\hypersetup{pdftitle={Local-rate passage with collective mean tracking},pdfauthor={Analytical increment for independent review}}
\begin{document}\maketitle
\begin{abstract}
We prove a nonlinear resident-passage estimate in the full leading
population. A fixed strong core is compared with a characteristic in
its own population field. The resident collective mode is uniformly
contracting for every strong mass between zero and one, while the core
spread has leading growth rate $\alpha(1-N)$ rather than
$\lambda(1-N)/2$. A coupled first-exit argument controls both the mean
and the spread, with an explicit Riccati denominator and no additional
logarithmic loss. The remaining strong population is retained through
its first absolute moment relative to the core reference; weak feedback
is retained through $N+|V|+K_w$. These integrated forcing budgets are
hypotheses, not initialized conclusions. A subsequent bounded interval
preserves the local power $N_0^{-\alpha/\lambda}$. We state the resulting
joint-parameter exponent test with every seed and shape $q$-power
visible. Neither Q2 retention, initialized entry, weak angular
condensation, nor finite-noise transfer is proved here.
\end{abstract}

\section{Scope and exact leading equations}\label{inc:AN03:local:v1:contract}
The source is the defined leading system in the supplied checkpoint,
\texttt{pa:leading}, \texttt{pa:masslaws}, \texttt{pa:overlap},
\texttt{pa:resident}, and \texttt{pa:spectra}. The pair cancellation is
also recorded in AN03 two-sided passage v1; its mixed-order proof is
repaired explicitly in v2. We restate all equations needed here. This is
not a uniform reduction theorem for the original positive-noise flow.

Fix $0<\rho<1$, $\lambda=\rho^2$, and
\[
 K(z)=\phi(z)+z\Phi(z),\qquad
 F(t,r)=\phi(\min(t,r))+r\Phi(\min(t,r)).
\]
Let $\mu_s,\mu_w$ be the strong and weak measures, with total moments
\[
 M=\int\dd\mu_s,\quad N=\int\dd\mu_w,\quad
 Y=\int y\dd\mu_s,\quad V=\int v\dd\mu_w,\quad
 K_w=\int K(u)\dd\mu_w.
\]
In particular $V$ is not a normalized mean. The strong field is
\begin{align}
 2f_x(x,y)&=-(1-M)x+\rho\phi(\rho x)-\rho\mathcal S(x)
       +\lambda[V+(1-N)y]\Phi(\rho x),\label{inc:AN03:local:v1:field}\\
 2f_y(x,y)&=-(1-M)y-Y-K_w+\rho(1-N)K(\rho x),\nonumber\\
 \mathcal S(x)&=\int F(\rho x,\rho\widetilde x)\dd\mu_s,\nonumber\\
 M'&=M(1-M),\qquad N'=\lambda N(1-N).\label{inc:AN03:local:v1:mass}
\end{align}
We consider characteristic solutions on finite intervals with bounded
support on compact subintervals, $0<M\le1$, $0<N\le1/2$, and fixed
family labels. Their normalized label weights are constant. All tail
statements below concern labels of this same leading system. A physical
population outside the noise-scaled charts needs an additional field
error budget; it is not incorporated by merely calling it a tail.

Let $\eta$ be the normalized full strong-label law. Fix a measurable
core $C$ with $\eta(C)>0$ at the initial time. No later relabeling is
allowed. Let $(c,d)$ be any characteristic of \eqref{inc:AN03:local:v1:field}
in the same evolving field; it need not carry mass. Define
\begin{align}
 D_\beta&=|x_\beta-c|+|y_\beta-d|,\qquad
 r=\esssup_{\beta\in C}D_\beta,\nonumber\\
 T&=\int_{C^c}D_\beta\dd\eta(\beta),\qquad
 \mathfrak f=N+|V|+K_w+T.\label{inc:AN03:local:v1:forcing}
\end{align}
Thus $\int D\dd\eta\le r+T$. The tail budget is a moment, not just a
mass fraction. A small mass at large scaled displacement need not have
small $T$.

Let $a=a_\rho$ be the unique positive solution of $a=\rho K(\rho a)$,
and put
\begin{align}
 P&=\Phi(\rho a),\qquad \alpha=\lambda P/2,\nonumber\\
 B&=a+1,\qquad C_0=4(a+3),\nonumber\\
 \gamma&=\frac{\alpha(1/2-\alpha)}{\alpha+1/2}>0,
 \qquad g=\gamma/2,\qquad
 e_*=\min\{1,\gamma/(2C_0)\}.\label{inc:AN03:local:v1:constants}
\end{align}
The symbol $e=\sqrt{(c-a)^2+(d-a)^2}$ denotes reference mean-direction
error, not the training residual.

\section{A local all-order shape inequality}\label{inc:AN03:local:v1:shape}
\begin{lemma}[Local coefficient with the full donor field]
\label{inc:AN03:local:v1:localineq}
While $e\le1$ and $r\le1$, the core radius obeys
\begin{equation}\label{inc:AN03:local:v1:rineq}
 D^+r\le[\alpha(1-N)+C_0(e+r+\mathfrak f)]r.
\end{equation}
No ordering between a core label and the reference is required.
\end{lemma}
\begin{proof}
Set $\varepsilon_M=1-M$ and define
\[
 \Pi(z)=\int(\widetilde x-z)_+\dd\mu_s,\quad
 \Lambda_-(z)=\int(z-\widetilde x)_+\dd\mu_s,\quad
 X_s=M^{-1}\int x\dd\mu_s.
\]
Then $\Pi-\Lambda_-=M(X_s-z)$. Also, by the explicit overlap derivative,
\begin{align*}
 2\partial_xf_x(z,d)
  &=-\varepsilon_M+\rho^3[V+(1-N)d-z-\Pi(z)]\phi(\rho z),\\
 2\partial_yf_x(z,d)&=2\partial_xf_y(z,d)
          =\lambda(1-N)\Phi(\rho z).
\end{align*}
Let $h=|x-c|$, $k=|y-d|$, $P(z)=\Phi(\rho z)$, and
\[
 U_c=V+(1-N)d-c,\qquad U_c^-=U_c+M(c-X_s),\qquad
 c_\rho=\rho^3\phi(0)/2.
\]
For $x>c$, subtraction and absolute values give, with
$\overline P_+=h^{-1}\int_c^{c+h}P(z)\dd z$,
\begin{align*}
 2D^+h&\le-\varepsilon_Mh+\lambda(1-N)kP(c+h)
                    +2c_\rho(U_c)_+h,\\
 2D^+k&\le-\varepsilon_Mk+\lambda(1-N)h\overline P_+.
\end{align*}
For $x<c$, use the donor cancellation before estimating. With
$\overline P_-=h^{-1}\int_{c-h}^cP(z)\dd z$ and $P_-=P(c-h)$,
\begin{align*}
 2D^+h&\le-\varepsilon_Mh+\lambda(1-N)kP_-
  +\lambda\varepsilon_Mh(\overline P_--P_-)
  +2c_\rho(U_c^-)_+h,\\
 2D^+k&\le-\varepsilon_Mk+\lambda(1-N)h\overline P_-.
\end{align*}
Both displays hold for either sign of $y-d$. At a zero coordinate
difference, the absolute value of its cross velocity gives the same
inequality. At full coincidence, uniqueness preserves coincidence.
Since $h\le r$, all displayed gate factors are at most $P(c+r)$.
The lower-side extra term is at most $\lambda\varepsilon_Mh$.
Consequently every core label satisfies
\[
 D^+D\le\left[-\frac{(1-\lambda)\varepsilon_M}{2}
   +\frac\lambda2(1-N)P(c+r)
   +c_\rho\max(0,U_c,U_c^-)\right]D.
\]
The Gaussian Lipschitz bound gives
$P(c+r)\le P+\rho\phi(0)(e+r)$. Furthermore,
\[
 \max(0,U_c,U_c^-)
 \le |V|+\sqrt2e+BN+r+T,
\]
since $|d|\le B$ and $|c-X_s|\le\int D\dd\eta\le r+T$.
Discard the negative strong-mass term and use $\lambda<1$, $\rho<1$,
$\phi(0)<1/2$; the constant $C_0$ in
\eqref{inc:AN03:local:v1:constants} exceeds all the resulting
coefficients. This proves the common labelwise bound with coefficient
$\alpha(1-N)+C_0(e+r+\mathfrak f)$.
On compact time intervals bounded velocities make all $D_\beta$
uniformly Lipschitz. Integrating their common scalar inequality and
taking the essential supremum gives \eqref{inc:AN03:local:v1:rineq}
in integral form and almost everywhere. The upper Dini interpretation
follows by continuity of the common coefficient; measurable forcing
may instead be kept in integral form.
\end{proof}

\section{Collective mean tracking, not an assumed mean budget}
\label{inc:AN03:local:v1:mean}
Collapse the entire strong mass to $(c,d)$, suppress weak feedback,
and write its field as $G_M$:
\begin{align}
 2(G_M)_1(c,d)&=-(1-M)c+\rho(1-M)\phi(\rho c)
                 +\lambda(d-Mc)\Phi(\rho c),\nonumber\\
 2(G_M)_2(c,d)&=-d+\rho K(\rho c).\label{inc:AN03:local:v1:coherentfield}
\end{align}
This is an analytical comparison field. The actual reference continues
to evolve in the full population field, not in $G_M$.

\begin{lemma}[Uniform stable mode and exact perturbation budget]
\label{inc:AN03:local:v1:meanineq}
For $e\le1$, the actual reference obeys
\begin{equation}\label{inc:AN03:local:v1:emean}
 D^+e\le-\gamma e+C_0e^2+C_0(r+\mathfrak f).
\end{equation}
Therefore, while $e\le e_*$,
\begin{align}
 e(t)&\le e_0e^{-g(t-t_0)}+
 C_0\int_{t_0}^te^{-g(t-s)}[r(s)+\mathfrak f(s)]\dd s,
 \label{inc:AN03:local:v1:convolution}\\
 \int_{t_0}^te&\le\frac{e_0}{g}
       +\frac{C_0}{g}\int_{t_0}^t(r+\mathfrak f).
 \label{inc:AN03:local:v1:intmean}
\end{align}
\end{lemma}
\begin{proof}
The point $(a,a)$ is stationary for $G_M$ at every $0\le M\le1$.
Its derivative is the symmetric matrix
\[
 J_M=\begin{pmatrix}-(1-M)/2-M\alpha&\alpha\\
                     \alpha&-1/2\end{pmatrix}
 =J_1-\begin{pmatrix}(1-M)(1/2-\alpha)&0\\0&0\end{pmatrix}.
\]
The positive matrix $-J_1$ has determinant
$\alpha(1/2-\alpha)$ and trace $\alpha+1/2$. Its smallest eigenvalue
is at least determinant divided by trace. Hence
$z^TJ_Mz\le-\gamma|z|^2$, uniformly in $M$.

For completeness, on $|c-a|^2+|d-a|^2\le1$ the only nonzero second
partials are bounded using
\begin{align*}
 2\partial_{cc}(G_M)_1
 &=-\rho^3[(M+1)+\rho^2c(d-c)]\phi(\rho c),\\
 2\partial_{cd}(G_M)_1&=\rho^3\phi(\rho c),\qquad
 2\partial_{cc}(G_M)_2=\rho^3\phi(\rho c).
\end{align*}
With $|c|,|d|\le B$ and $|d-c|\le2$, Taylor's formula on the segment
from $(a,a)$ bounds the vector remainder by $C_0e^2$.

The exact discrepancy $E=(c',d')-G_M(c,d)$ satisfies
\begin{align*}
 2E_1&=-\rho[\mathcal S(c)-MK(\rho c)]
                       +\lambda[V-Nd]\Phi(\rho c),\\
 2E_2&=Md-Y-K_w-\rho N K(\rho c).
\end{align*}
Since $0\le\partial_rF(t,r)\le1$,
\[
 |\mathcal S(c)-MK(\rho c)|\le\rho\int|x-c|\dd\mu_s.
\]
Also $|Y-Md|\le\int|y-d|\dd\mu_s$, $M\le1$, and
$K(\rho c)\le\phi(0)+\rho B$. Thus
\[
 |E|_2\le |E|_1\le\tfrac12(r+T)+\tfrac12(|V|+K_w)
                    +(B+1/2)N\le C_0(r+\mathfrak f).
\]
The Euclidean norm derivative, including its upper derivative at $e=0$,
proves \eqref{inc:AN03:local:v1:emean}. For $e\le e_*$, the quadratic
term is at most $\gamma e/2$. Scalar comparison yields
\eqref{inc:AN03:local:v1:convolution}; Tonelli and integration of the
exponential kernel yield \eqref{inc:AN03:local:v1:intmean}.
\end{proof}

\section{Coupled nonlinear local passage}\label{inc:AN03:local:v1:closure}
Fix $n_b\in(0,1/2]$, $0<N_0=N(t_0)<n_b$, and let $t_b$ be the unique
logistic time with $N(t_b)=n_b$. It is finite. Put
\[
 K_0=C_0(1+C_0/g),\qquad a_b=\alpha(1-n_b)>0.
\]

\begin{theorem}[Local-rate passage with proved mean tracking]
\label{inc:AN03:local:v1:main}
Assume the leading solution and reference exist on $[t_0,t_b]$ as in
Section 1, and the full weak/tail forcing satisfies
\[
 \int_{t_0}^{t_b}\mathfrak f\le F_b.
\]
Set
\begin{align}
 L_b&=\exp(C_0e_0/g+K_0F_b)(n_b/N_0)^{\alpha/\lambda},\nonumber\\
 \mathcal A&=r_0L_b,\qquad e_0=e(t_0),\quad r_0=r(t_0).
 \label{inc:AN03:local:v1:smallparameter}
\end{align}
If
\begin{align}
 e_0&\le e_*/4,\qquad F_b\le e_*/(4C_0),\nonumber\\
 \mathcal A&\le
 \min\left\{\frac14,\frac{a_b}{2K_0},
                      \frac{g e_*}{8C_0}\right\},
 \label{inc:AN03:local:v1:guards}
\end{align}
then throughout the interval
\begin{equation}\label{inc:AN03:local:v1:result}
 e\le3e_*/4,\qquad r\le2\mathcal A\le1/2,
 \qquad
 r(t_b)\le2e^{C_0e_0/g+K_0F_b}(n_b/N_0)^{\alpha/\lambda}r_0.
\end{equation}
The mean also satisfies the convolution bound
\eqref{inc:AN03:local:v1:convolution}, and
\begin{equation}\label{inc:AN03:local:v1:intresult}
 \int_{t_0}^{t_b}r\le\frac{2\mathcal A}{a_b},\qquad
 \int_{t_0}^{t_b}e\le\frac{e_0}{g}
              +\frac{C_0}{g}\left(F_b+\frac{2\mathcal A}{a_b}\right).
\end{equation}
In particular no lower bound on the actual growth of $r$ is needed.
\end{theorem}
\begin{proof}
Work until the first exit of $e<e_*$ or $r<1$. Integrate
\eqref{inc:AN03:local:v1:rineq} and use
\eqref{inc:AN03:local:v1:intmean}. Since
$\lambda\int_{t_0}^t(1-N)=\log(N(t)/N_0)$, this gives
\[
 r(t)\le r_0L(t)\exp\left(K_0\int_{t_0}^tr\right),\qquad
 L(t)=e^{C_0e_0/g}(N(t)/N_0)^{\alpha/\lambda}
                      e^{K_0\int_{t_0}^t\mathfrak f}.
\]
The case $r_0=0$ follows directly from same-field uniqueness. Otherwise,
differentiate $\exp(-K_0\int r)$ and integrate to obtain
\[
 r(t)\le\frac{r_0L(t)}{Q(t)},\qquad
 Q(t)=1-K_0r_0\int_{t_0}^tL(s)\dd s
\]
while $Q>0$. This argument is a scalar integral comparison, not an
assumption that the radius increases. The comparison multiplier obeys
\[
 L'=[\alpha(1-N)+K_0\mathfrak f]L\ge a_bL,
 \quad
 \int_{t_0}^tL\le[L(t)-L(t_0)]/a_b\le L_b/a_b.
\]
Consequently \eqref{inc:AN03:local:v1:guards} gives $Q\ge1/2$ and
$r\le2\mathcal A$, and also $\int r\le2\mathcal A/a_b$.
The convolution bound yields
\[
 e(t)\le e_0+\frac{2C_0\mathcal A}{g}+C_0F_b\le3e_*/4.
\]
Both first-exit guards improve strictly, so continuation proves the
claims on all of $[t_0,t_b]$. Finally integrate the same convolution to
obtain \eqref{inc:AN03:local:v1:intresult}.
\end{proof}

\paragraph{One way to supply the small forcing budget.}
If $|V|+K_w\le C_wN+w$, $w\ge0$, then
\[
 F_b\le\frac{1+C_w}{\lambda}\log\frac1{1-n_b}
                       +\int_{t_0}^{t_b}(w+T).
\]
Thus a fixed small $n_b$ and a small integrated weighted tail/weak
remainder supply the theorem's hypothesis. The latter estimates are not
proved here. A large initial weak chart radius cannot be replaced by a
bounded $C_w$ without an argument; an integral budget may still be useful.

\paragraph{Which mean is controlled.}
The core barycentre differs from $(c,d)$ in $\ell^1$ by at most $r$;
the full strong barycentre differs by at most $r+T$. The theorem therefore
controls the core mean through $e+r$, and the full mean through $e+r+T$.
An integral bound on $T$ does not itself give a pointwise full-mean bound.
Core radius about the core barycentre is at most $2r$.

\section{Bounded late propagation and explicit joint powers}
\label{inc:AN03:local:v1:powers}
\begin{proposition}[No further initialization power after fixed $n_b$]
\label{inc:AN03:local:v1:late}
Suppose, in addition, that the leading solution continues from $N=n_b$
to $N=1/2$, and all strong and weak coordinates and the reference are
bounded by $R\ge1$ on this remaining interval. Then
\[
 r(t_h)\le C_{\rm post}r(t_b),\qquad
 C_{\rm post}=\exp\left[
 \left(\frac\lambda2+5c_\rho R\right)
 \frac1\lambda\log\frac{1-n_b}{n_b}\right].
\]
This factor is independent of $N_0$.
\end{proposition}
\begin{proof}
The global pair bound in the proof of Lemma
\ref{inc:AN03:local:v1:localineq} has coefficient at most
$\lambda/2+c_\rho\max(0,U_c,U_c^-)$ after dropping strong damping.
Here $|V|\le NR\le R$, $|c|,|d|,|X_s|\le R$, so $|U_c|\le3R$
and $|U_c^-|\le5R$. The logistic time from $n_b$ to half-mass is
$\lambda^{-1}\log[(1-n_b)/n_b]$. Integrate the common labelwise bound.
\end{proof}

\begin{corollary}[Conditional joint-parameter exponent test]
\label{inc:AN03:local:v1:joint}
Suppose the hypotheses and forcing bounds of Theorem
\ref{inc:AN03:local:v1:main}, and where needed Proposition
\ref{inc:AN03:local:v1:late}, hold uniformly in a parameter family.
At the same chart-valid entry clock suppose
\[
 r_0\le C_s q^{-s_{\rm sh}}S_0^{\beta_{\rm sh}},\qquad
 N_0\ge c_w q^{s_{\rm seed}}S_0^{\gamma_{\rm seed}},
\]
with positive $C_s,c_w$ independent of $q,S_0$. Then the smallness
parameter and final core radius are bounded by a constant times
\begin{equation}\label{inc:AN03:local:v1:qpower}
 q^{-s_{\rm sh}-s_{\rm seed}\alpha/\lambda}
 S_0^{\beta_{\rm sh}-\gamma_{\rm seed}\alpha/\lambda}.
\end{equation}
For $S_0=q^\kappa$, a sufficient strict power condition is
\[
 \kappa(\beta_{\rm sh}-\gamma_{\rm seed}\alpha/\lambda)
 -s_{\rm sh}-s_{\rm seed}\alpha/\lambda>0.
\]
This does not prove any of the displayed entry laws.
\end{corollary}
\begin{proof}
Insert the two inequalities into $r_0(n_b/N_0)^{\alpha/\lambda}$.
The other factors are uniformly bounded by the hypotheses. A strict
positive $q$-power tends to zero and eventually satisfies the radius
smallness condition. The mean and forcing hypotheses remain separate.
\end{proof}

\paragraph{Candidate matching laws, explicitly conditional.}
The review proposes
\[
 \beta_{\rm sh}=\tfrac12-\alpha,\quad \gamma_{\rm seed}=1-\lambda,\quad
 s_{\rm sh}=\frac{1-2\alpha}{1-\lambda},\quad s_{\rm seed}=0.
\]
Under those \emph{unproved common-clock matching laws}, the power in
\eqref{inc:AN03:local:v1:qpower} is
\begin{equation}\label{inc:AN03:local:v1:candidate}
 E_{\rm loc}(\rho,\kappa)
 =\frac\kappa2(1-P)-\frac{1-\lambda P}{1-\lambda}.
\end{equation}
This is the local nonlinear exponent, not a global-rate exponent.

There is an elementary nonempty algebraic window on the reviewed ratio
interval. The rational root bound in AN03 repaired v2, also proved below,
gives $P<153/250$ for $\rho\in[13/20,7/10]$. At $\kappa=15/2$,
\eqref{inc:AN03:local:v1:candidate} decreases as either $P$ or $\lambda$
increases over this box. Hence
\begin{equation}\label{inc:AN03:local:v1:rationalwindow}
 E_{\rm loc}>
 \frac{15}{4}\frac{97}{250}
 -\frac{1-(49/100)(153/250)}{1-49/100}
 =\frac{4193}{51000}>0.
\end{equation}
For the bound on $P$, set $z=\rho a$, so $z=\lambda K(z)$.
Because $K(z)\le2/5+z/2+z^2/5$ for $z\ge0$ and
$(49/100)[2/5+7/50+49/3125]<7/25$, monotonicity of
$z-\lambda K(z)$ gives $z<7/25$. Thus
$P<1/2+(2/5)(7/25)=153/250$.

The same $\kappa$ lies strictly between the target-only crossing and
fixed-angle round-trip timing powers:
\[
 \frac2\lambda<\frac{15}{2}<\frac{2}{\lambda(1-\lambda)}
 \qquad(\lambda\in[169/400,49/100]).
\]
The right inequality follows from $\lambda(1-\lambda)\le1/4$;
the left from $2/(169/400)<15/2$. These facts identify a useful
\emph{conditional matching target}, not an initialized admissible
region for Theorem A. In particular the weak seed is not yet in a fixed
weak chart at the strong-learning clock in that proposed transit regime.
A later common entry clock must be proved and all its costs counted.

For reference, under the same candidate powers the algebraic condition
that $s_{\rm sh}/[(1-P)/2]$ be below $2/[\lambda(1-\lambda)]$ is exactly
\[
 (1+\lambda)P<1.
\]
Indeed multiply the comparison by positive denominators and factor
$1-\lambda$. This is a compatibility test for specified candidate laws,
not a phase boundary of the original population.

\section{Remaining initialized obligations}\label{inc:AN03:local:v1:residue}
The new theorem removes the assumption of an already controlled mean
trajectory during the local resident interval: its mean control is proved
simultaneously with the radius bound. It does not remove the required
entry mean/shape estimates or the weighted weak/tail forcing budget.
No estimate of tail mass alone supplies $\int T$.

To use the result for AN06 witnesses still requires: a common chart-valid
entry state; a correctly matched weak seed; weak angular relaxation and
selection of the desired coherent branch; the later bounded continuation;
and original finite-noise, gate-strip, and outside-population estimates.
The normalization of time remains the population's own weak half-mass
clock. Before chart capture, a coordinate-energy law cannot be silently
identified with the closed leading logistic mass law.

The statement that a radius has an upper differential growth bound does
not imply that the radius actually grows exponentially. Our proof bounds
its time integral using a growing comparison multiplier and its Riccati
denominator, so no such lower-growth assumption is used.

\begin{thebibliography}{9}\small\raggedright
\setlength{\itemsep}{2pt}\setlength{\parskip}{0pt}
\bibitem{P}\path{THEOREM_A_ANALYTICAL_PROGRESS.tex}, supplied checkpoint:
\texttt{pa:leading}, \texttt{pa:masslaws}, \texttt{pa:overlap},
\texttt{pa:resident}, \texttt{pa:spectra}.
\bibitem{A3}\path{AN03_two_sided_passage_and_Q2_retention_v1.tex} and
its separately repaired v2. The new local mean--shape closure is not
attributed to those sources.
\bibitem{R}\path{Pasted markdown.md}, user's review supplied in the
current turn. Reported data and matching powers are motivation, not proof
premises.
\bibitem{A2}\path{AN02_Q2_profiles_and_localized_tracking_v1.tex},
companion target-only and exact conditional tracking increment.
\end{thebibliography}
\end{document}
